\documentclass[a4paper,12pt,twoside,openright,titlepage]{article}

%Additional packages
\input{styles/packages}

% Listing style
\input{styles/lstset}

% Uncomment for production
% \syntaxonly

% Style
\pagestyle{headings}

% Define document
\author{D. Leeuw}
\title{Python Caesar Cipher}
\date{\today\\
1.0.0 \\
\vfill
\raggedright
\copyright\ 2025-2026 Dennis Leeuw\\
\input{styles/licentie-titlepage}}


\begin{document}
\selectlanguage{dutch}

\maketitle

%%%%%%%%%%%%%%%%%%%
%%% Introductie %%%
%%%%%%%%%%%%%%%%%%%

%B% \frontmatter
\section{Over dit Document}
\subsection{Leerdoelen}
\input{src/leerdoelen-builtin}
\subsection{Voorkennis}
\input{src/voorkennis-builtin-wb}

%%%%%%%%%%%%%%%%%
%%% De inhoud %%%
%%%%%%%%%%%%%%%%%
%%% \tableofcontents

%B% \mainmatter

\section{Caesar Cipher}
\input{src/caesar_cipher}

\subsection{Opdracht}
\input{src/opdracht-caesar_cipher}

\subsection{Tip: ASCII-tabel}
\input{src/ascii-table-intro}
\input{src/ascii-table}
\input{src/ascii-char2int}

 \subsection{Andere vertaaltabellen}
 \input{src/codetabellen}

\subsection{Tip: Python conversie functies}
\input{src/builtin-char2int}


%%%%%%%%%%%%%%%%%%%%%
%%% Index and End %%%
%%%%%%%%%%%%%%%%%%%%%
%\backmatter
\printindex
\end{document}

%%% Last line %%%
